% chktex-file 44
\newpage
\section{Methodisches Vorgehen}\label{Methodisches Vorgehen}

\subsection{Design Science Research nach Österle}
Die Entwicklung der Artefakte folgt dem iterativem Design Science Research Erkenntnisprozess nach Österle mit den idealtypischen Phasen\footcite[S. dazu][o. S.]{oesterle}:

\begin{figure}[H]
    \centering
    \includegraphics[width=0.9\textwidth]{oesterle}
    \caption{Design Science Research Erkenntnisprozess nach Österle}\cite[Quelle:~][S. 6]{oesterle}
\end{figure}

Diese Methode wurde gewählt, da sie es erlaubt (und auch fordert), dass die gebauten Prototypen aus dem Labor in ihrer realen Umgebung getestet werden. Es gibt bei diesem Projekt ein große Menge an Faktoren, die berücksichtigt werden müssen: Wetter, Licht, Strom \& WLAN, Temperatur \& Verdunsten der Lockflüssigkeit um nur ein paar zu nennen. Durch die iterativen Evaluationsphasen können die Prototypen kontinuierlich weiterentwickelt werden um sich Schrittweise einem perfekten Artefakt zu nähern. Ebenso erlaubt die Methode bei neuem Erkenntnissen aus der Evaluationsphase, wie beispielsweise, dass die Lockflüssigkeit keine Hornissen anlockt, wieder in den Analyseteile zurückzuspringen um ein Substitutionsprodukt zu suchen.

\subsection{Analyse}

\subsubsection{Problemstellung \& Motivation\label{sec:Problemstellung}}
Der Ursprung der Forschungsfrage kommt aus dem Interesse, die eigenen Bienen-Völker schützen zu wollen. Dieses Interesse wurde aber nach erster Recherche-Arbeit und Fachgesprächen um die Interessenvertretung aller Imker und Landwirte erweitert, da die entstehenden Schäden in Fauna und Wirtschaft bereits jetzt merkbar sind bzw. in Bayern zukünftig sein werden. Für die Grundlagenforschung wurde daher die FOM MyEBSCO Datenbank und Google Scholar nach öffentlich zugänglicher Literatur mit Kombinationen der Schlagworte \enquote{Asiatische Hornisse, asian hornet, Vespa velutina, Vespa velutina nigrithorax, Nestlokalisierung, Hornet nest location, Cluster Algorithmus, Cluster algorithm, Feldversuch, Feldexperiment, Detection} durchsucht. Außerdem wurden gängige Fachzeitschriften der Imkerszene und die Veröffentlichungen von staatlichen Institutionen sowie Interessensverbänden gesichtet, die zur Einschätzung der aktuellen Lage der Verbreitung und zu bewährten Methoden der Bekämpfung in der Praxis wertvolle Informationen lieferten. 

\subsubsection{Aktueller Forschungsstand}\label{sec:Forschungsstand}
Einen methodischen Ansatz, Nester ausfindig zu machen, liefert ein Forscherteam aus Spanien, die in einem Zwei-Phasen-Plan die Hornissen anlocken und markieren, um über die Rückflugrichtung zum Nest, sowie die Dauer des Flugs zurück zur Lock-Station Rückschlüsse auf die Nestposition zu ziehen. Falls das nicht sofort zum gewünschten Erfolg führt, können in einer zweiten Phase und durch weitere Standorte mittels Triangulierung die Nestpositionen genauer eingegrenzt werden\footcite[S. dazu][o. S.]{method}.

Versuche der Landesanstalt für Bienenkunde der Universität Hohenheim im Jahr 2025 mit speziellen Fallen und Lockstoffen die Verbreitung einzudämmen und zu dokumentieren, wurden als unzulässig eingestuft, da mangels ausreichender Selektion viele heimische (geschützte) Insekten in den Fallen verendeten\footcite[S. dazu][o. S.]{unihohenheim}.

Einen weniger destruktiven und technischen Ansatz wählten Forschende aus England, Portugal und Belgien bei ihren Versuchen den individuellen Ton der Flügelschläge einer asiatischen Hornisse mit Mikrofonen zu erfassen und durch maschinelles Lernen von anderen Insekten zu separieren\footcite[S. dazu][o. S.]{flugsound}. Ob und wie sich damit die Nester lokalisieren lassen, ist jedoch noch unklar. Deutlich zielführender scheint die Methode der Transponder und Radio-Telemetrie, die bereits bei anderen Tierarten seit Jahren erfolgreich zur Verfolgung eingesetzt wird. Dies stellt jedoch aufgrund der geringen Größe der Hornissen eine große Herausforderung dar und erfordert eine teure technische Grundausrüstung. Zusätzlich ermöglicht eine maximale Reichweite von 100 bis 150 Meter der Transponder nur einen begrenzten Anwendungsbereich und erfordert einen hohen zeitlichen Aufwand\footcite[S. dazu][o. S.]{radio-telemetry1}\footcite[S. dazu][o. S.]{radio-telemetry2}. 

Ein Forscherteam der University of Exter und dem Universitätsklinikum Tübingen trainierte mit einem Test-Datensatz ein \ac{CNN} namens VespAI, mit welchem die europäische und asiatische Hornisse über einen Raspberry Pi Computer und eine Kamera visuell zuverlässig erkannt und Alarmmeldungen verschickt werden können\footcite[S. dazu][o. S.]{vespai}. Eine Nestlokalisierung ist aber auch hierüber nicht möglich.

\subsubsection{Forschungslücke \& Anforderungen}
Es ist von größter Wichtigkeit eine Verwechslung mit der besonders schützenswerten heimischen Hornisse zu verhindern\footcite[Vgl.][S. 2]{aktionsplan}. Da 96 \% der Imker in Deutschland ihre Imkerei als Hobby betreiben\footcite[Vgl.][o. S.]{statista}, sind sowohl finanzielle Mittel als auch zeitliche Ressourcen für eine gezielte Bekämpfung der invasiven Hornissenart nur begrenzt vorhanden. Während die staatlichen Institute vermehrt auf die Mitarbeit der Bevölkerung bauen, um mit speziellen Dochttöpfen und Lockstoffen die Nester ausfindig zu machen\footcite[S. dazu][S. 1]{dochttoepfe}, gibt es aus der Wissenschaft mehrere technische Ansätze, die Ausbreitung der asiatischen Hornisse einzudämmen. 

Jedoch haben alle vorgestellten Ansätze das Problem, dass sie entweder sehr zeitaufwändig sind (Dochttöpfe, methodische Suche und Radio-Telemetrie), zu invasiv für andere Insekten (Fallen), oder die Kosten für das Equipment (Mikrofone oder Radio-Telemetrie) sehr hoch sind. Vor Allem aber sind es Einzelmaßnahmen, welche nicht auf Dauer oder in der Fläche skalieren und genau hier besteht eine Forschungslücke, welche mit dieser Arbeit geschlossen werden soll.

\subsubsection{Standortwahl}
Eine geeignete Standort für die Stationen zu finden ist nicht einfach, da diese möglichst Frei aus allen Richtungen angeflogen werden sollten doch auch Strom und ein WLAN-Verbindung benötigen. Am Wichtigsten jedoch, dass es tatsächlich auch \ac{aH} in der Gegend gibt. Da diese sich Stand 2026 erst den westlichen Grenzgebieten von Bayern ausgebreitet, werden die Stationen \enquote{PI-3} und \enquote{PI-4} in Rheinland-Pfalz von einem Mitglied des Imkervereins Bad-Dürkheim betreut. Die Stationen \enquote{PI-1} und \enquote{PI-2} werden im Münchner Westen erprobt. Hier gibt es zwar nur die europäische Hornisse, dafür können so leichter neue Releases getestet und bei Problemen schnell reagiert werden.

\subsection{Entwurf}

\subsubsection{Entwicklung des Artefakts Lock-Station}
\paragraph{Konstruktion der Lock-Station}

Bei der Konstuktion der Lock-Stationen wurde besonders auf geringe Kosten, Robustheit und Sicherheit wert gelegt, da die Stationen ein halbes Jahr dem Wetter trotzen müssen und einen 220V Stromanschluss benötigen. Die Technik, bstehend aus einem Raspberry Pi und einem Camera Modul sind in einem selbst-entwickelten und 3D-gedrucktem Gehäuse sicher geschützt unter einem Lampenschirm, welcher wiederum auf einer lakierten Holzplatte und einem umgedrehten Getränkekasten montiert ist. Letzteres soll dem Konstrukt Stabilität geben und einen sicheren und Wettergeschützen Aufbewahrungsort für die Lockflüssigkeit und Stromversorgung bieten.

\begin{figure}[H]
    \begin{subfigure}[c]{0.31\textwidth}
        \includegraphics[width=1\textwidth, angle = -90]{Lockstation1.jpeg}
        \subcaption{Station München1}
    \end{subfigure}
\begin{subfigure}[c]{0.31\textwidth}
        \includegraphics[width=1\textwidth, angle = -90]{Lockstation2.jpeg}
        \subcaption{Technik in 3D-Case}
    \end{subfigure}
\begin{subfigure}[c]{0.31\textwidth}
        \includegraphics[width=1\textwidth, angle = -90]{Lockstation3.jpeg}
        \subcaption{Kamera in 3D-Case}
    \end{subfigure}\\
    \centering{Quelle: Eigene Darstellung}
\end{figure}

Die Lockflüssigkeit entstammt dem Rezept\footcite[o. S.]{dochttoepfe} des \ac{LWG} und besteht aus 1/3 Weißwein, 1/3 dunkles Bier und 1/3 Himbeersirup und wird über eine Flasche und einem Docht von unten durch die Holzplatte in den Sichtbereich der Kamera geführt und mit einem Gummi nach oben gedrückt.

Beim Aufstellen muss darauf geachtet werden, dass die Kamera mit Hilfe des aufgedruckten Pfeils auf dem Case und einem Kompass nach Norden ausgerichtet wird und die Stromkabel wasserfest in Plastiktüten innerhalb des Getränkekasten verstaut wird.

\paragraph{Programmierung der Software}

Der Python Quellcode, welcher auf dem Raspberry-Pi einer Lockstation läuft, wird über Github verwaltet und kann so komfortabel über einfache \ac{SSH} Befehle oder der GUI \enquote{Raspberry Pi Connect} auf den Stationen aktuell gehalten werden. Die Software genügt den Prinzipien moderner Softwareentwicklung wie Kapselung \& Modularität und steht jedem unter der GPL-3.0 Lizenz über die Url \url{https://github.com/KarstenMS/hornet-radar/tree/main} zur freien Verfügung. Bei den Entwicklungsarbeiten hat anfangs ChatGPT\footcite[S. dazu][o. S.]{chatGPT} unterstützt, wurde aber bei zunehmender Komplexität von Claude Sonet der Firma Anthropic\footcite[S. dazu][o. S.]{claude} abgelöst.

Zum Einsatz kommen neben dem im nächsten Absatz beschriebenen VespAI KI Modell noch weiter Softwarepakete wie \enquote{Torch, Torchvision, OpenCV} als Paketabhängigkeit von VespAI, zur Bildverarbeitung oder um mit Background-Substraction Algorithmen wie \ac{MOG2} die Bewegungen von Insekten zu erkennen und zu isolieren (tracken)\footcite[Vgl.][o. S.]{mog2}.

Beim Aufruf ohne weitere Parameter startet die Software mit einem Debug Video feed, um die Ausrichtung zu kontrollieren und Informationen, wie \ac{FPS}, Konfidenz-Grenzwert zu sehen und auch das aktive Tracking und Klassifizierung live zu erleben. Im laufenden Betrieb muss das Programm als Hintergrund-Dienst ausgeführt und der Debug-Mode deaktiviert, um die Performance zu erhöhen. 

Sobald ein Bewegung erkannt wird, versucht ein Tracking Algorithmus das Objekt zu verfolgen und die Positionen in einem \gls{Dictionary} zu speichern. Bei stabilem Tracking wird anschließend über Vesp-AI eine Klassifizierung ausgeführt und wenn der Konfidenzintervall über dem Grenzwert liegt ein Beweisbild abgelegt und das Objekt weiter bis zu seinem Abflug verfolgt. Bei Verlust des Objekts werden aus den ersten und letzten Tracking Positionen ein An- und Abflugvektor berechnet und als zusammen mit Beweisbild, Verweildauer und weiteren Metadaten über eine Post-Web-Request in eine supabase Datenbank geschrieben. Zusätzlich werden die Events als .json in einem lokalen Ordner abgelegt und nach 6 Monaten automatisiert gelöscht.

Der Code bietet dem Anwender über eine config.py Konfigurationsdatei viele Parameter selbst zu bestimmen.
So können beispielsweise Kamera-Auflösung und Format, Distanz zum Lockstoff, Konfidenz-Grenzwerte und GPS-Standortdaten variabel vergeben werden - aber auch tiefgreifende Einstellungen wie der verwendete Tracking-Algorithmus, Farbanpassungen, Tracking-Box-Grenzen, Anzahl der Klassifizierungsversuche u.v.m.
Ebenfalls kann der Anwender über die config.py verschiedene Kamerasysteme und Optimierungen hinsichtlich Performance selbst einstellen. Eine hohe \ac{FPS} ist erstrebenswert um die schnell fliegenden Insekten \enquote{einfangen} zu können. Hier zeigen sich deutliche Unterschiede zwischen den Raspberry Pi Versionen. Ein Raspberry Pi 4 ist deutlich träger als die neuste Generation 5 und auf einem Raspberry Pi Nano konnte der Code gar nicht ausgeführt werden, da die Leistung der CPU ungenügend war.

\paragraph{Vesp-AI Objekterkennung}
Das Herzstück der Stationen ist die VespAI auf Basis von \ac{YOLO}v5 Objekterkennung. Das britische Forscherteam hat dazu das Modell mit drei Datasets, bestehend aus insgesamt 5449 Bildern von Hornissen (asiatisch und europäisch) und anderen Insekten trainiert und validiert\footcite[Vgl.][S. 2]{vespai}. Bei jeder Klassifizierung übergibt das Modell einen Konfidenzintervall-Wert, der angibt wie sicher sich das Modell bei seiner Vorhersage einer asiatische oder europäische Hornisse ist.

\ac{YOLO}v5 beruht auf dem Technikstand von 2021 und ist nicht mehr zeitgemäß. Auch die Abhängigkeiten zu älterer Python Softwarepaketen hat anfänglich Probleme bereitet. Auf Anfrage beim Entwicklungsteam wurde zwar ein neues und besseres Modell in Aussicht gestellt, jedoch wurden weitere Nachfragen ignoriert.

\paragraph{Precision vor Recall}
Damit in der späteren Cluster-Analyse für potentielle Nester keine False-Positives wie beispielsweise ziellose Fliegen und Wespen nicht die Ergebnisse verfälschen, ist eine hohe Precision bei der Objekterkennung von Vorteil. Die Precision gibt an, wie viele Objekte, die als Hornissen klassifiziert wurde, auch tatsächlich Hornissen waren. Zwar wäre auch ein hoher Recall wünschenswert, also wie viele tatsächliche Hornissen auch als Hornissen erkannt wurden, allerdings fallen verpasste Hornissen (und damit weniger Events) im späteren Clustering nicht so sehr ins Gewicht, wie irrtümlich generierte Events von beispielsweise einer Fliege.

Die Entwickler von Vesp-AI geben an, eine durchschnittliche Precision von \geq~0,99 und einen durschnittlichen Recall von > 0,93 über alle Kamera Systeme hinweg zu erreichen, betonen aber die Wichtigkeit einer optimal kalibrierten Kamera\footcite[Vgl.][S. 5]{vespai}.


\subsubsection{Entwicklung des Artefakts Website}
Für das erste Artefakt wurde zuerst eine Webhosting-Paket des Anbieters \enquote{lima-city} gebucht, um eine Domain mit SSL-Zertifikat und Wordpress-Installation zu erhalten. Dies erlaubt sowohl die technische Umsetzung des Clusterings via Javascript, als auch die Darstellung und Dokumentation der Entwicklung der Prototypen des Artefakts Lock-Station. Die Webseite ist über die UrL \url{https://hornet-radar.com} erreichbar. Da das Forschungsthema von internationalem Interesse ist, wird die Webseite sowohl als Deutsch als auch auf Englisch dargestellt.

\paragraph{Programmierung der interaktiven Karte mit Cluster-Algorithmus}
Die wichtigste Funktion der Webseite ist die interactive Karte von OpenStreetmap und dem Cluster-Algorithmus DB Scan. Bei Aufruf der Unterseite werden die Events der letzten 14 Tage aus supabase geladen und die Lock-Stationen als Punkte auf der Karte dargestellt. Gibt es ausreichend viele Events, die den Standard-Werten \enquote{\varepsilon ~Nachbarschaft} und \enquote{minPts} entsprechen, werden diese Events über eine Javascript DBSCAN-Worker Funktion zu Clustern zusammengefasst und als farbliche Korridore mit Statistik-Informationen dargestellt. Besucher und \enquote{Nest-Finder} können über Schiebe-Regler die Parameter für \enquote{DBScan, Zeitraum und Reichweite} variieren und so die Triangulation der potentiellen Nestposition beeinflussen. So kann je nach Datenmenge dünne und genaue Korridore mit einer hohen Kohärenz oder dicke und ungenaue Korridore mit einer niedrigen Kohärenz erzeugt werden.
Dort wo sich die Korridore schneiden, wäre eine Nestersuche am Erfolgsversprechenden. 

\begin{figure}[H]
    \centering
    \includegraphics[width=0.9\textwidth]{karte}
    \caption{Karte mit Cluster-Korridoren}
    \centering{Quelle: \url{https://hornet-radar.com/de/map}}
\end{figure}
   
\paragraph{Programmierung der Detailansicht für Events von Lock-Stationen}
In der Unterseite Detailansicht werden alle Events der jeweiligen Lock-Stationen zeitlich absteigend sortiert ausgegeben. Jede Zeile steht für ein Event, mit Datum, Uhrzeit, Konfidenzintervall, An- und Abflugwinkel, sowie einem Beweisbild, damit jeder Stadionbetreiber die Basis der Cluster-Korridore nachvollziehen kann.

\paragraph{Programmierung des Polardiagramms zur Darstellung der An- und Abflugrichtung}
Das Polardiagram entstand aus der Not heraus, die Cluster-Analyse besser verstehen zu wollen. Zwar kann auf der Detailseite für jede Station und Event der An- und Abflugwinkel nachverfolgt werden, doch ist es deutlich komfortabler die Messpunkte gesammelt in einem Polardiagramm zu sehen. Für die generelle Funktionalität hat diese Seite jedoch keine Relevanz.

\begin{figure}[H]
    \centering
    \includegraphics[width=0.9\textwidth]{polar}
    \caption{Polardiagramm mit An- und Abflugwinkeln der Events}
    \centering{Quelle: \url{https://hornet-radar.com/de/polar}}
\end{figure}

\subsubsection{Kostenaufstellung für Lock-Station und Website}

Damit die Lock-Stationen bei Imkern flächendeckend eingesetzt werden, muss der Kosten-Nutzen-Faktor stimmen. Da ausschließlich OpenSource Software zum Einsatz kommt, belaufen sich die laufenden Kosten lediglich auf das Web-Hosting Paket und die Doamin-Registrierung für die Webseite, was jedoch mit 3,50 € im Monat vernachlässigbar ist. Anders würde es bei einer Kommerzialisierung diese Lösung oder bei einer großen vertikalen Skalierung aussehen, da hier die OpenSource Lizenzen meist technisch begrenzt sind oder eine Gebühr verlangt wird. 

Bei den Fix-Kosten sind vor allem die Raspberry Pi die großen Kostenträger wie die Tabelle \nameref{Kosten-Tabelle} zeigt. Die Speicherkosten-Explosion der letzten Monate hat zudem die Preise noch weiter in die Höhe getrieben. Hat ein Raspberry Pi 5 B mit 4 GB RAM im Januar 2026 noch durchschnittlich 71,90 € gekostet sind es im August bereits 105,90 €.\footcite[Vgl.][o. S.]{idealo}

\begin{table}[H]
    \caption{Konfiguration \& Netto-Kosten der Lockstation}\label{Kosten-Tabelle}
    \small                          % etwas kleinere Schrift hilft enorm
    \setlength{\tymin}{1.5cm}       % Mindestbreite für tabulary-Spalten
    \begin{tabulary}{\textwidth}{@{}L>{\centering\arraybackslash}p{2.5cm}>{\centering\arraybackslash}p{2.5cm}>{\centering\arraybackslash}p{2.5cm}>{\centering\arraybackslash}p{2.5cm}@{}}
    %\begin{tabulary}{\textwidth}{@{}LLLLL@{}}
        \toprule
        \textbf{Element \& Preis} & \textbf{Station M{\"u}nchen 1} & \textbf{Station M{\"u}nchen 2} & \textbf{Station Freinsheim 3} & \textbf{Station Freinsheim 4} \\
        \midrule
        Raspberry Pi                & v. 5 B 4GB RAM        & v. 5 B 4GB RAM    & v. 5 B 4GB RAM        & v. 4 B 4GB RAM    \\[4pt]
                                    & 86,90 €               & 86,90 €           & 86,90 €               & 60,90 €           \\[4pt]
        Kamera Modul                & Module 3              & Module 3          & Module 3              & Pi AI Camera      \\[4pt]
                                    & 28,90 €               & 28,90 €           & 28,90 €               & 83,99 €           \\[4pt] 
        USB-C Netzteil 27W          & 12,40 €               & 12,40 €           & 12,40 €               & 12,40 €           \\[4pt]
        3D Case                     & 5 €                   & 5 €               & 5 €                   & 5 €               \\[4pt]
        Holzplatte \& Kleinteile    & 15 €                  & 15 €              & 15 €                  & 15 €              \\[4pt]  
        \bottomrule
        Gesamtkosten                & 148,20 €              & 148,20 €          & 148,20 €              & 177,29 €          \\[4pt]

    \end{tabulary}
    \\    \vspace{2pt}
    {\centering Quelle: Eigene Darstellung}
\end{table}


\subsection{Evaluation}

\subsubsection{Evaluationskriterien}

Für eine vergleichbare und objektive Bewertung der Lock-Stationen und Software werden folgende Kriterien definiert, welche nach jeder Iteration erhoben und für die Anpassungen der nächsten Iteration analysiert werden.

\begin{itemize}
    \item[$\circ$] \textbf{Precision} 
    Sie gibt an, wie viele der gesamten visuellen Erkennung von asiatischen und europäischen Hornisse durch VespAI korrekt waren und dadurch für eine brauchbaren Cluster-Analyse taugen.
    Alles unter 50 \% wäre schlecht, 50-80 \% mittel und alles darüber als gut zu betrachten. Dieser Wert lässt dich durch den Konfidenzintervall des KI Modells beeinflussen
    
    \item[$\circ$] \textbf{Anzahl} 
    Die generelle Anzahl von Sichtungen Hornissen in relation zu allen erzeugten Events um genügend Daten für eine zuverlässige Clustering-Bildung zu erhalten. Die Anzahl erlaubt außerdem den Standort der Lock-Station zu beurteilen.

    \item[$\circ$] \textbf{Invasivität} 
    Gezählt werden die Anzahl verendete Insekten
    \item[$\circ$] \textbf{FPS} 
    Die \ac{FPS} gibt Aufschluss über die Performance der Systeme. Eine hohe \ac{FPS} vermindert das Risiko, dass das Tracking abbricht dadurch viele kurze und unnütze Events entstehen. Es werden einmal die FPS im Ruhezustand gemessen und einmal mit einem aktiven Tracking eines Insekts.
    \item[$\circ$] \textbf{Qualität Resultate}
    Eine subjektive Bewertung der übertragenen Werte. Taugen die Beweisbilder und ergeben die gelieferten An- und Abflugwinkel und andere Metadaten einen Sinn. 
    \item[$\circ$] \textbf{Qualität Station} 
    Eine generelle Bewertung der Stationen hinsichtlich Verschleiss und Wetter-Beständigkeit.   
\end{itemize}

\subsubsection{Iteration Mai}

\begin{itemize}
    \item \textbf{Ergebnisse}

    Die erste Testphase von einem Monat ergab bisher kein einziges Event mit einer asiatischen oder europäischen Hornisse, dafür einige Ameisen und Wespen an den Stationen. Es ist noch früh in der Saison, in welcher die Hornissen bevorzugt nach Proteinquellen suchen um ihre Brut zu versorgen. Kurioser Weise gibt es ein Event, bei welcher eine Wespe mit 86 \% Konfidenz als \ac{aH} klassifiziert wurde, die daneben sitzende asiatische Hornisse jedoch nicht.
    
    \begin{figure}[H]
    \centering
    \includegraphics[width=0.5\textwidth]{false-detection.png}
    \caption{Falsche Klassifizierung einer Wespe}
    \centering{Quelle: Eigene Darstellung}
    \end{figure}

    Dass das System jedoch prinzipiell funktioniert zeigt, dass einige Events von Wespen und Fliegen generiert wurden, welche auf Grund des bewusst niedrig gehaltenen Konfidenzintervalls jedoch fälschlicherweise als Hornissen klassifiziert wurden. Für einen Bewertung des Cluster-Algorithmus sind diese Events jedoch unbrauchbar.
    In der Lockflüssigkeit wurden an zwei Stationen eine tote Wespen gefunden, die an dem Docht vorbei in die Lockflüssigkeit gefallen und gestorben sind. 
    Bei der Performance Schnitt der ältere Raspberry Pi 4 deutlich schlechter ab, was angesichts der CPU-Unterschiede nicht verwunderlich ist. Jedoch fallen die \ac{FPS} besonders bei aktivem Tracking von Insekten deutlich ein und können bei laufender VespAI Klassifizierung kurzzeitig bis auf einen \ac{FPS} absacken, was ein kontinuierliches Tracking unmöglich macht.
    Auch sind Unterschiede in Qualität und Farbe der Bilder festzustellen, welche auf Grund künstlicher Lichtquellen im Winter nicht direkt aufgefallen sind. 

    Das größte Problem ist jedoch, die teils große Distanz zu den WLan-Routern wodurch es immer wieder zu Verbindungsabbrüchen und Tageslangen Offline-Zeiten kommt. Dies ist besonders kritisch, da die Lock-Stationen nur dann Events generieren können, wenn sie online sind und die Bilder an den Server übertragen werden können. Es erschwert aber auch die Administration und das Verteilen neuer Software-Releases.


    \begin{table}[H]
        \caption{Ergebnisse Mai}
        \small                          % etwas kleinere Schrift hilft enorm
        \setlength{\tymin}{1.5cm}       % Mindestbreite für tabulary-Spalten
        \begin{tabulary}{\textwidth}{@{}L>{\centering\arraybackslash}p{2.3cm}>{\centering\arraybackslash}p{2.3cm}>{\centering\arraybackslash}p{2.3cm}>{\centering\arraybackslash}p{2.3cm}@{}}
        %\begin{tabulary}{\textwidth}{@{}LLLLL@{}}
            \toprule
            \textbf{Evaluation} & \textbf{Station M{\"u}nchen 1} & \textbf{Station M{\"u}nchen 2} & \textbf{Station Freinsheim 3} & \textbf{Station Freinsheim 4} \\
            \midrule
            Precision                               & 0 \%          & 0 \%          & 0 \%          & 0 \%          \\[4pt]
            \# Hornissen / Events                   & 0 / 6         & 0 / 10        & 0 / 22        & 0 / 29        \\[4pt]
            Invasivität                             & 0             & 1             & 0             & 1             \\[4pt]
            \varnothing~FPS ruhend / tracking       & 7,0 / 5,3     & 7,7 / 3,7     & 6,7 / 6,6     & 5,4 / 3,3     \\[4pt]  
            Qualität Resultate                      & Mittel        & Schlecht      & Mittel        & Mittel        \\[4pt]
            Qualität Station                        & Gut           & Gut           & Gut           & Gut           \\[4pt]
            \bottomrule
        \end{tabulary}
        \\    \vspace{2pt}
        {\centering Quelle: Eigene Darstellung}
    \end{table}


    \item \textbf{Anpassungen für den Juni Zyklus}
    \begin{itemize}
        \item[$\circ$] \textbf{Hardware:}\\ 
        Der Lockstoff-Docht wurde mit einem Knoten verdickt, damit keine Ameisen und Wespen in die Lock-Flüssigkeit fallen können. 
        \item[$\circ$] \textbf{Software:}\\ 
        Ein Cron-Job startet nun die Raspberry Pi jeden Morgen um 6 Uhr einmal durch, damit wieder eine Verbindung zum WLAN hergestellt werden kann. 
        Außerdem wird der Konfidenz-Grenzwert für die Klassifizierung von 75 \% auf 90 \% erhöht, um die vielen False-Positives zu reduzieren. 
        Zusätzlich wurden statische teile der Konfiguration wie die GPS-Lokation und Pi-ID in eine lokalen config\_local.py Datei ausgelagert, welche nicht durch Änderungen am Code permanent von Git überschrieben wird. 
        Um die Qualität und die \ac{FPS} zu erhöhen wird der Auto-Fokus der Camera Module 3 deaktiviert. Das AI-Camera Module verfügt über einen manuellen Zoom, welcher nur geändert werden kann, wenn das Modul ausgebaut wird. Da dies ein langwieriges und aufwendiges Verfahren ist, disqualifiziert sich die AI-Kamera für eine weitere Verwendung und wird bei nächster Gelegenheit durch ein Camera Module 3 ersetzt.
        Des weiteren wurde ein Bug bei der Konvertierung von BGR zu RGB Bildformat behoben, wodurch VespAI und die Webseite keine Bilder mehr mit Blaustich bekommen:
        \begin{figure}[H]
            \begin{subfigure}[c]{0.35\textwidth}
                \includegraphics[width=1\textwidth]{bluepicture}
                \subcaption{BGR-Format}
            \end{subfigure}
            \begin{subfigure}[c]{0.335\textwidth}
                \includegraphics[width=1\textwidth]{originalpicture}
                \subcaption{RGB-Format}
            \end{subfigure}
        \caption{Vergleich BGR- / RGB-Format beim Beispiel einer Wespe}
        {\centering Quelle: Eigene Darstellung}
        \end{figure}
        \end{itemize} 
\end{itemize}


\subsubsection{Iteration Juni}

\begin{itemize}
    \item \textbf{Ergebnisse}
    Die Verfügbarkeits-Änderungen haben den gewünschten Erfolg liefern können. Die Stationen sind dauerhaft erreichbar und erzeugen über den gesamten Zeitraum hinweg Events. Ebenfalls sind durch den erhöhten Konfidenzintervall die vielen False-Positives reduziert worden. Hornissen wurden immer noch nicht gesichtet, doch gibt es immer noch viele Fliegen und Wespen, welche fälschlicherweise als Hornissen klassifiziert werden. 
    Die Performance hat sich leider nicht merklich verbessert und besonders bei viel Bewegung (durch Fliegen, Ameeisen, etc.) kommen die Raspberry Pi an ihre Grenzen und aus dem niedrigen 1-3 FPS Bereich kaum mehr raus. Dafür hat das Abschalten des Auto-Fokus die Qualität der Bilder deutlich verbessern können, so dass schnelle Bewegungen die zentralen Elemente mehr \enquote{verwaschen}. 
    Einige der verwendeten Haushalts-Gummi, welche die Flasche mit der Lockflüssigkeit nach Oben drücken sollen, haben die Temperaturschwankungen nicht überstanden und sind gerissen. In gegen ursprünglicher Annahme sind die Flaschen mit der Lockflüssigkeit noch nahezu voll und könnten so noch einige Monate ohne Nachfüllen auskommen. 
    
   
    \begin{table}[H]
        \caption{Ergebnisse Juni}
        \small                          % etwas kleinere Schrift hilft enorm
        \setlength{\tymin}{1.5cm}       % Mindestbreite für tabulary-Spalten
        \begin{tabulary}{\textwidth}{@{}L>{\centering\arraybackslash}p{2.3cm}>{\centering\arraybackslash}p{2.3cm}>{\centering\arraybackslash}p{2.3cm}>{\centering\arraybackslash}p{2.3cm}@{}}
        %\begin{tabulary}{\textwidth}{@{}LLLLL@{}}
            \toprule
            \textbf{Evaluation} & \textbf{Station M{\"u}nchen 1} & \textbf{Station M{\"u}nchen 2} & \textbf{Station Freinsheim 3} & \textbf{Station Freinsheim 4} \\
            \midrule
            Precision                               & 0 \%          & 0 \%          & 0 \%          & 0 \%          \\[4pt]
            \# Hornissen / Events                   & 0 / 3         & 0 / 7         & 0 / 4         & 0 / 29        \\[4pt]
            Invasivität                             & 0             & 1             & 0             & 1             \\[4pt]
            \varnothing~ FPS ruhend / tracking      & 9,0 / 6,3     & 8,7 / 5,7     & 8,7 / 8,6     & 5,6 / 3,2     \\[4pt]  
            Qualität Resultate                      & Gut           & Gut           & Gut           & Mittel        \\[4pt]
            Qualität Station                        & Mittel        & Gut           & Mittel        & Mittel        \\[4pt]
            \bottomrule
        \end{tabulary}
        \\    \vspace{2pt}
        {\centering Quelle: Eigene Darstellung}
    \end{table}


    \item \textbf{Anpassungen für Juli Zyklus}
    \begin{itemize}
        \item[$\circ$] \textbf{Hardware:}\\ Die Gummis wurden durch widerstandsfähigere Kabelbinder ersetzt. Das macht das Nachfüllen der Lockflüssigkeit etwas aufwendiger, jedoch scheint dies nicht so häufig notwendig zu sein als ursprünglich angenommen.
        \item[$\circ$] \textbf{Software:}\\   Als Maßnahme für eine weiter verbesserte Performance wurde die automatische Belichtung durch feste Werte in der lokalen Konfigurationsdatei ersetzt. 
        Außerdem werden prozentuale Min- und Max-Werte für die Tracking Boxen eingeführt, damit auch bei wechselnder Auflösungen beispielsweise keine Ameisen getrackt werden und so die CPU-Last reduziert wird. Außerdem wird der Konfidenzintervall weiter auf 93 \% erhöht, um die Precision weiter zu erhöhen.
    \end{itemize}
\end{itemize}


\subsubsection{Iteration Juli}
\begin{itemize}
    \item \textbf{Ergebnisse}
    Auch in diesem Monat wurden keine asiatische Hornissen gesichtet. Ein Umfrage im lokalen Imker Verein bestätigt allerdings, dass es dieses Jahr in den frühen Monaten nur zu sehr vereinzelten Sichtungen kam (siehe Anhang \ref{Kommunikation}). 
    Erfreulicherweise die Stationen in München von europäischen Hornissen regelrecht belagert. Die Erkennungsrate von VespAI versagt hier jedoch komplett. Zur Kontrolle wurde ein Video über mehrere Sekunden aufgezeichnet und mit dem original Code von VespAI auf ihrer GitHub-Seite überprüft, jedoch mit ähnlichem Resultat:

    \begin{figure}[H]
    \begin{subfigure}[c]{0.30\textwidth}
        \includegraphics[width=1\textwidth, angle = -90]{Wrong-Detection3.jpeg}
        \subcaption{Falsche Klassifizierung M{\"u}nchen 1}
    \end{subfigure}
\begin{subfigure}[c]{0.30\textwidth}
        \includegraphics[width=1\textwidth, angle = -90]{Wrong-Detection4.jpeg}
        \subcaption{Falsche Klassifizierung M{\"u}nchen 2}
    \end{subfigure} \\

    \begin{subfigure}[c]{0.30\textwidth}
        \includegraphics[width=1\textwidth, angle = -90]{Wrong-Detection5.jpeg}
        \subcaption{Falsche Klassifizierung VespAI Code}
    \end{subfigure}
\begin{subfigure}[c]{0.30\textwidth}
        \includegraphics[width=1\textwidth, angle = -90]{Wrong-Detection6.jpeg}
        \subcaption{Falsche Klassifizierung VespAI Code}
    \end{subfigure} \\
    \centering{Quelle: Eigene Darstellung}
\end{figure}


    Ein weiteres großes Problem ist nach wie vor der Verlust des Trackings nach einer Klassifizierung auf Grund der geringen \ac{FPS} Rate. 
    Das deaktivieren der automatischen Belichtung hingegen hat zu stark überlichteten und sehr dunklen Bilder geführt. Da Hornissen Nachtaktiv sind und viele Events in den Morgen- und Abendstunden entstehen, wurde diese Einstellung nach einigen Tagen wieder auf automatisch zurückgestellt.
    \begin{table}[H]
        \caption{Ergebnisse-Juli}
        \small                          % etwas kleinere Schrift hilft enorm
        \setlength{\tymin}{1.5cm}       % Mindestbreite für tabulary-Spalten
        \begin{tabulary}{\textwidth}{@{}L>{\centering\arraybackslash}p{2.3cm}>{\centering\arraybackslash}p{2.3cm}>{\centering\arraybackslash}p{2.3cm}>{\centering\arraybackslash}p{2.3cm}@{}}
        %\begin{tabulary}{\textwidth}{@{}LLLLL@{}}
            \toprule
            \textbf{Evaluation} & \textbf{Station M{\"u}nchen 1} & \textbf{Station M{\"u}nchen 2} & \textbf{Station Freinsheim 3} & \textbf{Station Freinsheim 4} \\
            \midrule
            Precision                               & 0 \%          & 0 \%          & 0 \%          & 0 \%          \\[4pt]
            \# Hornissen / Events                   & 0 / 3         & 0 / 7         & 0 / 4         & 0 / 29        \\[4pt]
            Invasivität                             & 0             & 0             & 0             & 0             \\[4pt]
            \varnothing~ FPS ruhend / tracking      &               &           &      &     \\[4pt]  
            Qualität Resultate                      &         &       &         &         \\[4pt]
            Qualität Station                        &         &            & Mittel        &         \\[4pt]
            \bottomrule
        \end{tabulary}
        \\    \vspace{2pt}
        {\centering Quelle: Eigene Darstellung}
    \end{table}

    \item \textbf{Anpassungen für August-Oktober Zyklus}
    \begin{itemize}
        \item[$\circ$] \textbf{Hardware:}\\ 
        \item[$\circ$] \textbf{Software:}\\ 
    \end{itemize}
\end{itemize}

\iffalse{}

Implementierung Multi-Tracking, um die Erkennung von mehreren Hornissen gleichzeitig zu ermöglichen.
Deaktivierung Belichtung um FPS zu erhöhen und Tracking zuverlässiger zu machen
Belichtung aufgrund von Nacht Aktiven Hornissen und wenig Performance Gewinn wieder deaktiviert
Sortierung der Ergebnisse im Detailview und Pi-ID mit in die Ansicht
Yolo Detection in eigenen Worker Daemon-Thread und auf 2 Kerne begrenzen um unabhängig vom Tracking zu laufen < bisher bester Performance gewinn



\subsubsection{Iteration September / Oktober}











\subsection{Diffusion}




\fi